% !TEX program = pdflatex
% Demo Day pitch deck — 3 sequential sections:
%   Passerby  (1--6):   60-second hook, self-contained
%   Professor (7--23):  full story, self-contained
%   Extra     (24--30): deep-dive Q\&A, flip here from anywhere
% Reuses FDR presentation conventions.
\documentclass[10pt,aspectratio=169]{beamer}

% -- Theme ---------------------------------------------------------------
\usetheme{Madrid}
\usecolortheme{default}
\setbeamertemplate{navigation symbols}{}
\setbeamertemplate{footline}[frame number]

% -- Packages ------------------------------------------------------------
\usepackage{graphicx}
\usepackage{booktabs}
\usepackage{tabularx}
\usepackage{listings}
\lstset{basicstyle=\footnotesize\ttfamily,breaklines=true}
\usepackage{hyperref}
\usepackage{amssymb}

% -- Convenience macros --------------------------------------------------
\newcommand{\incfig}[2][0.8\textwidth]{%
  \IfFileExists{#2}{%
    \includegraphics[width=#1, height=0.5\textheight, keepaspectratio]{#2}%
  }{%
    \fbox{\makebox[#1][c]{\shortstack{\small FIGURE TBD\\\tiny\detokenize{#2}}}}%
  }%
}
\newcommand{\incfigw}[2][\textwidth]{\incfig[#1]{#2}}

% One-line takeaway at bottom of content slides --- presenter prompt, viewer gist.
\newcommand{\takeaway}[1]{%
  \vspace{0.25cm}
  \begin{center}\textit{\small #1}\end{center}
}

% ======================================================================
% HOW TO DUPLICATE A SLIDE
% ------------------------
% To make the professor path self-contained (so you can jump to it without
% losing the passerby), copy any passerby slide's frame block and paste it
% into the professor section.  For example:
%
%   1. Copy the entire \begin{frame}...\end{frame} block from slide 2.
%   2. Paste it at the start of the professor section (after the section
%      header comment) as a new slide.
%   3. Give the new slide a % -- Slide N: ... comment matching its position.
%
% Repeats currently removed to keep every slide visually unique for review.
% Add them back when you finalize the deck for presentation.
% ======================================================================
\title{TrayGuard --- Surgical Tray Training}
\subtitle{EC ENGR 180DW --- Final Design Review}
\author{Owen Fan \and Matthew Li \and Andy Ma \and Jacob Moore \and Qiyu Xin}
\date{June 2026}

\begin{document}

% ======================================================================
%  SECTION 1: PASSERBY PITCH (slides 1--6, ~60 sec)
% ======================================================================

% -- Slide 1: Title ----------------------------------------------------
\begin{frame}
  \titlepage
\end{frame}

% -- Slide 2: The Problem ------------------------------------------------
\begin{frame}{The Problem}
  \begin{columns}[T]
    \column{0.55\textwidth}
    \begin{itemize}
      \item[\textbf{--}] \textbf{9 defects per 100} surgical cases

      \item[\textbf{--}] \textbf{55\%} occur during tray assembly

      \item[\textbf{--}] \textbf{10 min} average delay per affected case

      \item[\textbf{--}] \textbf{\$6.75--9.42M} estimated annual cost per hospital

      \item[\textbf{--}] \textbf{44\%} of packaging errors: wrong specification
    \end{itemize}
    \vspace{0.4cm}
    \begin{alertblock}{}
      \centering
      \emph{``The total count was correct. The specific instrument was wrong.''}
    \end{alertblock}
    \column{0.45\textwidth}
    \incfig{figures/teammate_kelly_crile.jpg}
    \vspace{0.1cm}
    {\tiny Kelly vs Crile --- same family, different serration pattern}
  \end{columns}
\end{frame}

% -- Slide 3: Why Cameras Failed -----------------------------------------
\begin{frame}{Why Cameras Failed}
  \begin{columns}[T]
    \column{0.55\textwidth}
    \begin{itemize}
      \item[\textbf{--}] Seven staged experiments: brightness, background, layout occlusion

      \item[\textbf{--}] \textbf{Every real condition} degraded detection

      \item[\textbf{--}] Occlusion: \textbf{mAP50--95 dropped to 0.389}

      \item[\textbf{--}] The model learned \textbf{positions, not shapes}

      \item[\textbf{--}] Site-specific data needed per hospital --- not feasible for one project
    \end{itemize}
    \vspace{0.4cm}
    \begin{alertblock}{}
      \centering
      The data told us no.
    \end{alertblock}
    \column{0.45\textwidth}
    \incfig{figures/overlay_map50_95.png}
    {\tiny mAP50--95 degradation under occlusion}
  \end{columns}
\end{frame}

% -- Slide 4: TrayGuard ---------------------------------------------------
\begin{frame}{TrayGuard}
  \begin{columns}[T]
    \column{0.55\textwidth}
    \begin{itemize}
      \item[\textbf{--}] Training platform: \textbf{laptop + printer + cards}

      \item[\textbf{--}] Camera scans which cards are on the table

      \item[\textbf{--}] Scoring engine classifies every selection into \textbf{6 error categories}

      \item[\textbf{--}] Confidence capture: separates confidently wrong from fragile knowledge

      \item[\textbf{--}] No special equipment. No hospital deployment required.
    \end{itemize}
    \vspace{0.4cm}
    \begin{exampleblock}{}
      \centering
      The real bottleneck is training the people who build trays.
    \end{exampleblock}
    \column{0.45\textwidth}
    \incfig{figures/card_front_example.png}
    \vspace{0.15cm}
    \incfig{figures/card_back_example.png}
  \end{columns}
\end{frame}

% -- Slide 5: How It Works ------------------------------------------------
\begin{frame}{How It Works}
  \centering
  \incfigw[0.88\textwidth]{figures/learning_loop.png}
  \vspace{0.4cm}
  \begin{columns}[T]
    \column{0.33\textwidth}
    \centering
    \textbf{Pre / Post Test} \\[0.2cm]
    \small Parallel variants \\[0.1cm]
    \small Different photo views \\[0.1cm]
    \small No hints \\[0.1cm]
    \small Baseline + measured gain
    \column{0.33\textwidth}
    \centering
    \textbf{Study + Quiz} \\[0.2cm]
    \small Retrieval-first cards \\[0.1cm]
    \small Prompt-before-reveal \\[0.1cm]
    \small Immediate feedback \\[0.1cm]
    \small Confidence capture
    \column{0.34\textwidth}
    \centering
    \textbf{Practice Sort} \\[0.2cm]
    \small Physical cards on table \\[0.1cm]
    \small Scored against count sheet \\[0.1cm]
    \small 6 error categories \\[0.1cm]
    \small Weak-item review
  \end{columns}
\end{frame}

% -- Slide 6: What We Need ------------------------------------------------
\begin{frame}{What We Need}
  \centering
  \vspace{0.8cm}
  {\Large Prototype built. Study protocol designed.}\\[1.2cm]
  {\large\bfseries Looking for a hospital partner\\to pilot the first real tray module.}\\[1.2cm]
  \begin{block}{}
    \centering
    Open-source. Documented design handoff.\\[0.3cm]
    \small Contact the TrayGuard team through UCLA EC ENGR 180DW.
  \end{block}
  \vspace{0.6cm}
  {\footnotesize Direct material cost: \textbf{\$53}. Requires: printer + existing laptop or tablet.}
\end{frame}

% ======================================================================
%  SECTION 2: PROFESSOR PATH (slides 7--23, ~4--5 min)
%  Add repeats of slides 2--6 here to make this section self-contained.
%  See duplication instructions in the preamble comments.
% ======================================================================

% -- Slide 7: Error by the Numbers ----------------------------------------
\begin{frame}{Evidence: Error by the Numbers}
  \vspace{0.3cm}
  \begin{columns}[T]
    \column{0.32\textwidth}
    \begin{block}{\centering\small}
      \textbf{Alfred et al. (2021)} \\[0.2cm]
      3,900 defects \\ 41,799 cases \\[0.2cm]
      \textbf{9 per 100} \\[0.1cm]
      \textbf{55\% at assembly}
    \end{block}
    \column{0.32\textwidth}
    \begin{block}{\centering\small}
      \textbf{Nichol et al. (2024)} \\[0.2cm]
      236 errors \\ 147 cases \\[0.2cm]
      \textbf{10 min avg delay} \\[0.1cm]
      \textbf{\$6.75--9.42M / yr}
    \end{block}
    \column{0.32\textwidth}
    \begin{block}{\centering\small}
      \textbf{Zhu et al. (2019)} \\[0.2cm]
      398 errors \\ 33,839 packages \\[0.2cm]
      \textbf{44\% wrong spec} \\[0.1cm]
      Largest error category
    \end{block}
  \end{columns}
  \vspace{0.5cm}
  \begin{center}
    Three different methods converge on the same finding:\\[0.2cm]
    tray errors are recurring, measured, operational failures.
  \end{center}
  \takeaway{Three methods. One finding.}
\end{frame}

% -- Slide 8: Kelly vs Crile Deep Dive -----------------------------------
\begin{frame}{The Lookalike Problem: Kelly vs Crile}
  \begin{columns}[T]
    \column{0.48\textwidth}
    \centering
    \incfig{figures/teammate_kelly_crile.jpg}
    \vspace{0.1cm}
    {\tiny Kelly (left) vs Crile (right) --- serration pattern is the
    discriminating feature}
    \column{0.48\textwidth}
    \begin{itemize}
      \item[\textbf{--}] Same instrument family: both hemostatic clamps

      \item[\textbf{--}] \textbf{Crile:} serrations cover the full jaw

      \item[\textbf{--}] \textbf{Kelly:} serrations only cover distal half

      \item[\textbf{--}] Sizes vary: 6'', 7'', 8''

      \item[\textbf{--}] 44\% of packaging errors are wrong specification
    \end{itemize}
  \end{columns}
  \vspace{0.3cm}
  \takeaway{A few millimeters of serration --- invisible at a glance.}
\end{frame}

% -- Slide 9: Root Cause Fishbone -----------------------------------------
\begin{frame}{Root Cause Analysis}
  \centering
  \incfigw{figures/fishbone_rca.png}
  \vspace{0.2cm}
  \begin{center}
    \footnotesize Six categories: instrument design, knowledge, process,
    information, environment, policy.\\[0.1cm]
    \footnotesize Three structural drivers: SPD as cost center, no licensure,
    no OR→SPD information system.
  \end{center}
  \takeaway{Six causes. Three structural drivers.}
\end{frame}

% -- Slide 10: Broken Feedback Loop ----------------------------------------
\begin{frame}{The Broken Feedback Loop}
  \centering
  \incfigw{figures/feedback_loop.png}
\end{frame}

% -- Slide 11: 7 CV Experiments --------------------------------------------
\begin{frame}{Evidence: 7 Staged CV Experiments}
  \vspace{0.2cm}
  \footnotesize
  \begin{tabularx}{\textwidth}{l X c}
    \toprule
    Experiment & Condition & Impact \\
    \midrule
    Brightness (train dark, test bright) & 80 → 800 lumens & mAP degrades with gap \\
    Brightness (train bright, test dark) & 800 → 80 lumens & Same pattern, symmetric \\
    Background reflectivity & Matte vs reflective & Matte best, reflective worst \\
    Layout: separated → separated & Clean lab layout & mAP50 = 0.990 \\
    Layout: separated → overlay & Instruments touching & mAP50--95 = 0.389 \\
    Real instruments (kitchen proxy) & 6 classes, fixed positions & Per-class AP 0.000--0.862 \\
    Real instruments (position check) & Split-1 vs Split-2 & Classes survive different splits \\
    \bottomrule
  \end{tabularx}
  \vspace{0.3cm}
  \takeaway{Every real condition degraded detection.}
\end{frame}

% -- Slide 12: Occlusion Before / After ------------------------------------
\begin{frame}{Evidence: Occlusion --- Before and After}
  \begin{columns}[T]
    \column{0.48\textwidth}
    \centering
    \incfig{figures/separated_layout_example.jpg}
    \vspace{0.1cm}
    {\tiny Separated --- what you train on}
    \column{0.48\textwidth}
    \centering
    \incfig{figures/overlay_layout_example.jpg}
    \vspace{0.1cm}
    {\tiny Overlay --- what a real tray looks like}
  \end{columns}
  \vspace{0.3cm}
  \takeaway{You'd have to rearrange instruments to scan them.}
\end{frame}

% -- Slide 13: Position vs Shape -------------------------------------------
\begin{frame}{Evidence: Position vs Shape Learning}
  \begin{itemize}
    \item[\textbf{--}] Per-class AP ranged from \textbf{0.000 to 0.862} depending on data split

    \item[\textbf{--}] Different instrument classes survived in different splits

    \item[\textbf{--}] Only the \textbf{scalpel} transferred --- the most visually distinct

    \item[\textbf{--}] Each image had one of each class at a \textbf{fixed position} per session

    \item[\textbf{--}] Model learned \textbf{which class at which location}, not visual features

    \item[\textbf{--}] When position changed (different tray), the model failed
  \end{itemize}
  \vspace{0.2cm}
  \takeaway{The model memorized positions, not shapes.}
\end{frame}

% -- Slide 14: The Pivot: TrayGuard -------------------------------
\begin{frame}{The Pivot: TrayGuard}
  \begin{columns}[T]
    \column{0.55\textwidth}
    \begin{itemize}
      \item[\textbf{--}] Training platform: \textbf{laptop + printer + cards}

      \item[\textbf{--}] Camera scans which cards are on the table

      \item[\textbf{--}] Scoring engine classifies every selection into \textbf{6 error categories}

      \item[\textbf{--}] Confidence capture: separates confidently wrong from fragile knowledge

      \item[\textbf{--}] No special equipment. No hospital deployment required.
    \end{itemize}
    \vspace{0.4cm}
    \begin{exampleblock}{}
      \centering
      The real bottleneck is training the people who build trays.
    \end{exampleblock}
    \column{0.45\textwidth}
    \incfig{figures/card_front_example.png}
    \vspace{0.15cm}
    \incfig{figures/card_back_example.png}
  \end{columns}
\end{frame}

% -- Slide 15: Error Categories -------------------------------------------
\begin{frame}{Design: Error Categories}
  \begin{columns}[T]
    \column{0.55\textwidth}
    \begin{itemize}
      \item[\textbf{--}] \textbf{Correct} --- item present, correct count

      \item[\textbf{--}] \textbf{Missing} --- required item not selected

      \item[\textbf{--}] \textbf{Extra} --- selected item not required

      \item[\textbf{--}] \textbf{Wrong} --- wrong instrument family

      \item[\textbf{--}] \textbf{Misidentified} --- lookalike substitution

      \item[\textbf{--}] \textbf{Wrong Count} --- correct instrument, wrong quantity
    \end{itemize}
    \vspace{0.3cm}
    \begin{block}{}
      \centering
      \footnotesize Each category maps to published literature.\\
      \footnotesize When a learner improves, we know \emph{which type} of error they reduced.
    \end{block}
    \column{0.45\textwidth}
    \centering
    \textbf{Published evidence}
    \vspace{0.2cm}
    \begin{itemize}
      \item \textbf{Alfred et al.} --- missing, extra, wrong
      \item \textbf{Nichol et al.} --- wrong, delays
      \item \textbf{Zhu et al.} --- 44\% misidentified
    \end{itemize}
  \end{columns}
  \takeaway{Same categories hospitals already track.}
\end{frame}

% -- Slide 16: Confidence Capture -----------------------------------------
\begin{frame}{Design: Confidence Capture}
  \begin{columns}[T]
    \column{0.55\textwidth}
    \vspace{0.3cm}
    \begin{itemize}
      \item Confidence rated per item: \textbf{1--5 scale}

      \item Plus a \textbf{``Not sure''} flag

      \item \textbf{High-confidence error}: confidently wrong
      \begin{itemize}
        \item Hardest to correct --- learner doesn't know they need help
      \end{itemize}

      \item \textbf{Low-confidence correct}: fragile knowledge
      \begin{itemize}
        \item May not survive a week without review
      \end{itemize}

      \item Both are exported in CSV/JSON
    \end{itemize}
    \column{0.45\textwidth}
    \incfig{figures/confidence_mockup.png}
  \end{columns}
  \vspace{0.2cm}
  \takeaway{Confidently wrong is not the same as fragile knowledge.}
\end{frame}

% -- Slide 17: Scope Boundaries --------------------------------------------
\begin{frame}{Scope Boundaries}
  \vspace{0.3cm}
  \begin{columns}[T]
    \column{0.48\textwidth}
    \begin{block}{\centering What TrayGuard Does}
      \begin{itemize}
        \item Local tray familiarization
        \item Lookalike discrimination practice
        \item Simulated tray sort with error-category feedback
        \item Pre/post measurement with CSV/JSON export
      \end{itemize}
    \end{block}
    \column{0.48\textwidth}
    \begin{alertblock}{\centering What It Does NOT Do}
      \begin{itemize}
        \item Replace CRCST certification
        \item Inspect instrument damage or wear
        \item Fix staffing shortages or wages
        \item Close the OR→SPD feedback loop
      \end{itemize}
    \end{alertblock}
  \end{columns}
  \vspace{0.4cm}
  \takeaway{Claim bounded to: improvement on a simulated task in one session.}
\end{frame}

% -- Slide 18: Alternatives Considered -------------------------------------
\begin{frame}{Alternatives Considered}
  \vspace{0.2cm}
  \footnotesize
  \begin{tabularx}{\textwidth}{l X}
    \toprule
    Approach & Why not buildable by a student team \\
    \midrule
    Automated tray checking (CV) & Per-hospital datasets, cross-manufacturer
    validation, regulatory review \\
    Physical error-proofing (foam) & Requires instrument inventory, manufacturer
    relationships, per-tray tooling \\
    Workflow redesign (dual check) & Requires policy authority, institutional
    buy-in, staffing changes \\
    Decision-support kiosks & Requires deployment access, workflow integration,
    multi-stakeholder approval \\
    Tray rationalization & Requires hospital utilization data, surgeon
    preferences, multi-department coordination \\
    Competency tracking systems & Requires IT integration, HIPAA compliance,
    ongoing institutional commitment \\
    \textbf{Training} & \textbf{Buildable with cards + laptop. Validatable with
    pre/post study. Creates durable local content any future solution needs.} \\
    \bottomrule
  \end{tabularx}
  \vspace{0.3cm}
  \takeaway{Training was the only buildable path.}
\end{frame}

% -- Slide 19: How It Works (repeat of slide 5) ----------------------------
\begin{frame}{Design: How It Works}
  \centering
  \incfigw[0.88\textwidth]{figures/learning_loop.png}
  \vspace{0.4cm}
  \begin{columns}[T]
    \column{0.33\textwidth}
    \centering
    \textbf{Pre / Post Test} \\[0.2cm]
    \small Parallel variants \\[0.1cm]
    \small Different photo views \\[0.1cm]
    \small No hints \\[0.1cm]
    \small Baseline + measured gain
    \column{0.33\textwidth}
    \centering
    \textbf{Study + Quiz} \\[0.2cm]
    \small Retrieval-first cards \\[0.1cm]
    \small Prompt-before-reveal \\[0.1cm]
    \small Immediate feedback \\[0.1cm]
    \small Confidence capture
    \column{0.34\textwidth}
    \centering
    \textbf{Practice Sort} \\[0.2cm]
    \small Physical cards on table \\[0.1cm]
    \small Scored against count sheet \\[0.1cm]
    \small 6 error categories \\[0.1cm]
    \small Weak-item review
  \end{columns}
\end{frame}

% -- Slide 20: Study Ladder -----------------------------------------------
\begin{frame}{Evaluation: The Study Ladder}
  \centering
  \incfigw{figures/study_ladder.png}
  \vspace{0.3cm}
  \small
  \begin{tabular}{l l}
    \textbf{Study 0} & Marker reliability --- do the cards scan? \\
    \textbf{Study 1} & Novice baseline --- within-subject pre/post, 8--12 participants \\
    \textbf{Study 2} & 7-day retention --- does improvement persist? \\
    \textbf{Study 3} & Expert review --- SPD educator validates content \\
    \textbf{Study 4} & SPD workplace pilot --- test with real technicians \\
  \end{tabular}
  \takeaway{Each step validates the next.}
\end{frame}

% -- Slide 21: Falsification Criteria -------------------------------------
\begin{frame}{Evaluation: Falsification Criteria}
  \centering
  \vspace{0.8cm}
  {\large For the 7-day retention check (Study 2):}\\[1cm]
  \begin{enumerate}
    \item \textbf{Delayed accuracy retains at least half} of the immediate gain
    \vspace{0.6cm}
    \item \textbf{At least 5 percentage points above baseline}
    \vspace{0.6cm}
    \item \textbf{At least 70\% of participants return} for the follow-up
  \end{enumerate}
  \vspace{0.8cm}
  \begin{alertblock}{}
    \centering
    If any criterion fails: the claim reverts to \textbf{``immediate learning only.''}
  \end{alertblock}
  \takeaway{We specify upfront what would falsify our claim.}
\end{frame}

% -- Slide 22: What We Accomplished ----------------------------------------
\begin{frame}{What We Accomplished}
  \vspace{0.2cm}
  \begin{itemize}
    \item[\textbf{--}] \textbf{7 staged CV experiments} with documented negative result

    \item[\textbf{--}] \textbf{Root cause analysis} across 6 categories + 3 structural drivers

    \item[\textbf{--}] \textbf{Full system specification:} 7-step learning loop, scoring engine, confidence capture, CSV/JSON export

    \item[\textbf{--}] \textbf{Study protocol} with explicit falsification criteria

    \item[\textbf{--}] \textbf{Prototype web app} (FastAPI + vanilla JS) implementing the full learner loop

    \item[\textbf{--}] \textbf{Durable tray module schema} --- JSON files capturing local instrument knowledge
  \end{itemize}
  \vspace{0.3cm}
  \takeaway{A rigorous no on CV is a real finding.}
\end{frame}

% -- Slide 23: What We Need (repeat of slide 6) ----------------------------
\begin{frame}{What We Need}
  \centering
  \vspace{0.8cm}
  {\Large Prototype built. Study protocol designed.}\\[1.2cm]
  {\large\bfseries Looking for a hospital partner\\to pilot the first real tray module.}\\[1.2cm]
  \begin{block}{}
    \centering
    Open-source. Documented design handoff.\\[0.3cm]
    \small Contact the TrayGuard team through UCLA EC ENGR 180DW.
  \end{block}
  \vspace{0.6cm}
  {\footnotesize Direct material cost: \textbf{\$53}. Requires: printer + existing laptop or tablet.}
\end{frame}

% ======================================================================
%  SECTION 3: EXTRA Q\&A (slides 24--30)
% ======================================================================

% -- Slide 24: Not Just Flashcards -----------------------------------------
\begin{frame}{Q\&A: How Is This Different from Flashcards?}
  \vspace{0.2cm}
  \footnotesize
  \begin{tabularx}{\textwidth}{l X X}
    \toprule
    Feature & Flashcards & TrayGuard \\
    \midrule
    Assessment & Recall only & Physical tray sort against count sheet \\
    Feedback & Binary right/wrong & 6 error categories mapped to hospital tracking \\
    Context & Isolated instrument name & Distractors, lookalikes, required quantities \\
    Confidence & None & 1--5 scale + ``Not sure'' flag \\
    Output & None & CSV/JSON export with hashed learner ID \\
    \bottomrule
  \end{tabularx}
  \vspace{0.3cm}
  \takeaway{The applied task, not just memorization.}
\end{frame}

% -- Slide 25: Why AprilTags -----------------------------------------------
\begin{frame}{Q\&A: Why AprilTags?}
  \vspace{0.2cm}
  \begin{itemize}
    \item[\textbf{--}] \textbf{Faster detection} than QR codes --- fewer bits to decode

    \item[\textbf{--}] \textbf{Tolerates glare and partial occlusion} through built-in error correction

    \item[\textbf{--}] \textbf{Stable 6-DOF pose estimation} handles camera angle variation

    \item[\textbf{--}] \textbf{Manual fallback:} 3-digit ID on card back --- type it if the camera misses

    \item[\textbf{--}] Cards are re-printable paper. No real instruments needed.
  \end{itemize}
  \vspace{0.2cm}
  \takeaway{Designed for tabletop cameras, not phones.}
\end{frame}

% -- Slide 26: Anti-Memorization Design ------------------------------------
\begin{frame}{Q\&A: Anti-Memorization Design}
  \begin{columns}[T]
    \column{0.40\textwidth}
    \centering
    \incfig{figures/card_front_example.png}
    \vspace{0.1cm}
    {\tiny Assessment cards: neutral front --- no name, no hint}
    \column{0.60\textwidth}
    \begin{itemize}
      \item[\textbf{--}] Pre-test and post-test use \textbf{matched parallel variants}

      \item[\textbf{--}] Assessment cards show only the \textbf{AprilTag on the front}

      \item[\textbf{--}] Pre and post use \textbf{different photo views} (A / B)

      \item[\textbf{--}] Random seed ensures \textbf{unique layouts per session}
    \end{itemize}
  \end{columns}
  \vspace{0.2cm}
  \takeaway{You cannot pass by memorizing card positions.}
\end{frame}

% -- Slide 27: Customizable Tray Module ------------------------------------
\begin{frame}[fragile]{Q\&A: Customizable Tray Module}
  \begin{columns}[T]
    \column{0.52\textwidth}
    \begin{verbatim}
{
  "instrument_id": "hospital_kelly_6in",
  "display_name": "Clamp Kelly 6\"",
  "family": "Clamping/occluding",
  "aliases": ["KELLY 6\"",
              "Kelly clamp 6"],
  "lookalike_pairs": [
    { "expected": "kelly_6in",
      "selected": "crile_5_5",
      "feedback": "Check serration
                    pattern..." }
  ],
  "required_items": [
    { "instrument": "kelly_6in",
      "quantity": 2 }
  ]
}
    \end{verbatim}
    \column{0.48\textwidth}
    \begin{itemize}
      \item[\textbf{--}] Single editable JSON file per hospital tray

      \item[\textbf{--}] Instruments, local names, aliases, counts

      \item[\textbf{--}] Distractors, lookalike pairs, features

      \item[\textbf{--}] No code changes --- edit file, reprint cards

      \item[\textbf{--}] One module = \textbf{2--4 hours} SPD educator time

      \item[\textbf{--}] Reusable for every new hire, competency check
    \end{itemize}
  \end{columns}
  \vspace{0.15cm}
  \takeaway{Edit the file. Reprint the cards.}
\end{frame}

% -- Slide 28: Per-Hospital Data Burden ------------------------------------
\begin{frame}{Q\&A: Per-Hospital Data Burden}
  \vspace{0.2cm}
  \begin{itemize}
    \item[\textbf{--}] Each hospital uses \textbf{different manufacturers, makes, and models}

    \item[\textbf{--}] SPD lighting varies by \textbf{workstation, time of day, fixture condition}

    \item[\textbf{--}] Tray layouts differ by \textbf{procedure type and surgeon preference}

    \item[\textbf{--}] Cross-manufacturer validation needed --- accuracy drops across brands

    \item[\textbf{--}] Conservative estimate: \textbf{thousands of annotation hours per hospital}

    \item[\textbf{--}] For \textbf{\textasciitilde5,500} US hospitals: industry-scale infrastructure problem
  \end{itemize}
  \vspace{0.2cm}
  \takeaway{Not a one-project problem.}
\end{frame}

% -- Slide 29: Cost \& Next Steps ------------------------------------------
\begin{frame}{Q\&A: Cost \& Next Steps}
  \begin{columns}[T]
    \column{0.50\textwidth}
    \centering
    \textbf{Prototype Bill of Materials}
    \vspace{0.3cm}
    \begin{tabularx}{\columnwidth}{l r}
      \toprule
      Item & Cost \\
      \midrule
      Printed cards (15--16 instruments) & \$15--30 \\
      Tray surface / mat & \$5--10 \\
      Webcam (if laptop cam insufficient) & \$10--15 \\
      \midrule
      \textbf{Total} & \textbf{\$30--55} \\
      \bottomrule
    \end{tabularx}
    \vspace{0.4cm}
    {\footnotesize Assumes existing laptop or tablet.}
    \column{0.50\textwidth}
    \textbf{Priorities for the Next Team}
    \vspace{0.3cm}
    \begin{enumerate}
      \item Run novice baseline study

      \item Expert review of tray content

      \item 7-day retention check

      \item SPD workplace pilot

      \item Hospital partner for real module
    \end{enumerate}
  \end{columns}
  \vspace{0.3cm}
  \takeaway{\$53 in materials. Clinical validation first.}
\end{frame}

% -- Slide 30: End ---------------------------------------------------------
\begin{frame}
  \centering
  \vfill
  {\Huge TrayGuard}\\[1cm]
  {\large Thank you. Questions?}\\[1cm]
  {\small Contact the TrayGuard team through UCLA EC ENGR 180DW.}\\[0.5cm]
  {\footnotesize Open-source. Documented design handoff.}
  \vfill
\end{frame}

\end{document}
